\documentclass[11pt]{article}

\usepackage[margin=1.8cm]{geometry}
\usepackage[T1]{fontenc}
\usepackage{microtype}
\usepackage{enumitem}
\usepackage[hidelinks]{hyperref}
\usepackage{xcolor}
\usepackage{ragged2e}

\setlength{\parindent}{0pt}
\setlength{\parskip}{5pt}
\setlist[itemize]{leftmargin=1.25em,itemsep=2pt,topsep=2pt}
\pagestyle{empty}
\urlstyle{same}

\begin{document}

\begin{flushright}
\textbf{St\'{e}phane Ga\"{e}l R. Ekodeck}\\
Department of Computer Science, University of Yaound\'{e} I\\
LIRIMA, Team GRIMCAPE, Yaound\'{e}, Cameroon\\
\href{mailto:stephane-gael.ekodeck@facsciences-uy1.cm}
{stephane-gael.ekodeck@facsciences-uy1.cm}\\
ORCID: 0000-0002-8094-8832\\
June 5, 2026
\end{flushright}

\vspace{2pt}

The Editor-in-Chief\\
\textit{IEEE Transactions on Dependable and Secure Computing}

\textbf{Re: Regular Paper submission, ``FAIR-BETH: A Fair Thresholding and
Ablation Protocol for Behavioral Malicious-Process Detection under
Distribution Shift''}

Dear Editor-in-Chief,

We respectfully submit the above manuscript for consideration as a Regular
Paper in \textit{IEEE Transactions on Dependable and Secure Computing}
(TDSC). The paper addresses a concrete dependability problem in
security-machine-learning systems: a detector can rank malicious processes
well yet produce unreliable operational decisions when models receive
different threshold treatment or when validation and deployment distributions
differ.

\textbf{Scope and relevance to TDSC.}
The work concerns dependable evaluation of endpoint-security detectors under
class imbalance, scarce malicious validation data, and distribution shift.
It is directly relevant to TDSC because it studies when a learned security
component can and cannot support a reproducible alerting decision. The paper
connects model evaluation to system dependability through explicit controls
for test isolation, feature and label leakage, calibration, common operating
policies, uncertainty, early-detection limits, host/time separation, and
response boundaries.

\textbf{Scientific gap and contribution.}
Published security-ML guidance warns against data leakage, weak baselines, and
inappropriate metrics. However, these principles do not by themselves provide
an executable procedure for deciding whether a hybrid detector adds value
under the same operating constraints as a strong non-hybrid baseline.
FAIR-BETH closes this gap through:
\begin{itemize}
    \item an evaluation contract that fixes data access, feature fitting,
    calibration, threshold selection, baseline parity, and uncertainty
    reporting before official-test evaluation;
    \item three quantitative diagnostics---tabular dominance, score
    compression, and threshold sensitivity---that make fusion claims
    falsifiable;
    \item repeated cross-fitted threshold selection over all available BETH
    development positives, without synthetic enrichment;
    \item target-copy-free GRU/LSTM capacity ablations that bound the
    sequence-model conclusion; and
    \item independent protocol validation on MalBehavD-V1, supplemented by
    BETH prefix, calibration, attribution, cost, feature-stress,
    host-disjoint, and temporal audits.
\end{itemize}

\newpage

\textbf{Principal evidence.}
After process aggregation, the official BETH training and validation files
contain no malicious processes, whereas 121 of 198 official-test processes
are malicious. Under one common validation-FPR policy, RF-500 and TabRF obtain
F1 scores of 0.7241 and 0.7210, while MetaRF obtains 0.6756; the hybrid fusion
therefore provides no evidence of improvement. Fifty thresholds fixed by
repeated five-fold validation have median 0.006 and IQR 0.004--0.006. Yet the
locked median threshold converts a 5\% development-FPR objective into 32.47\%
test FPR, with F1 0.6964. This separates threshold sampling stability from
threshold transfer and is the paper's main dependability finding.

The ranking--decision distinction is independently visible in ExtraTrees,
which achieves the best BETH ranking (AUC 0.7973, AP 0.8572) but only F1
0.6180 at its locked operating point. Target-copy-free recurrent models remain
weak across 23\,666--241\,074 parameters (AUC 0.4160--0.4911), so the paper
limits this negative result to the tested objectives and representation. On
MalBehavD-V1, the same evaluation contract yields repeated-split median F1
0.9533, supporting protocol portability without claiming direct cross-schema
model transfer.

\textbf{Claim boundaries and reproducibility.}
The manuscript does not claim ransomware-family attribution, calibrated
posterior risk, direct BETH-to-MalBehavD model transfer, or readiness for
autonomous containment. These exclusions follow from the available labels,
schemas, test size, and early-detection results. The reproducibility package,
including executable scripts, generated CSV/JSON artifacts, the main
manuscript, and separate supplementary material, is publicly available at
\href{https://github.com/EkodeckStephane/fair-beth-malicious-process-detection}
{github.com/EkodeckStephane/fair-beth-malicious-process-detection}.

\textbf{Mandatory author declarations.}
\begin{itemize}
    \item \textbf{Originality and exclusive submission:} The manuscript is
    original, has not been published, is not under consideration elsewhere,
    and will not be submitted to another venue while under TDSC review.
    \item \textbf{Author approval:} All named authors have read and approved
    the manuscript, its authorship order, and the stated contributions.
    \item \textbf{Conflicts and funding:} The authors declare no competing
    interests. The research received no external funding.
    \item \textbf{Ethics and data:} The study involves no human participants,
    animals, or sensitive personal data. It uses publicly available
    cybersecurity datasets and reports derived experimental artifacts.
    \item \textbf{Supplementary material:} Detailed feature specifications,
    hyperparameters, threshold sensitivity, data/model cards, deployment
    boundaries, and audit figures are supplied as a separate file, consistent
    with TDSC submission requirements.
\end{itemize}

We believe the manuscript will interest TDSC readers working on dependable
security analytics, evaluation under distribution shift, endpoint detection,
and reproducible security-machine-learning methodology. Thank you for your
consideration.

\vspace{4pt}
Yours sincerely,

\vspace{15pt}

\textbf{St\'{e}phane Ga\"{e}l R. Ekodeck}\\
Corresponding Author\\
On behalf of Vivien Beyala Kamgang, Serge Alain Eb\'{e}l\'{e}, Ulrich Joel
Atchiengang Ngueyep, St\'{e}phane Willy Mossebo Tcheunteu, Damaris Belle M.
Fotso, Arthur Ulrich Ewane, and Julius Marcellin Nkenlifack

\end{document}
